\section{Two bounded variation quadrature estimates}\label{sec:quadrature}
We first isolate the deterministic geometric estimate used at every stage.
Let $U$ denote uniform probability on $[0,1]$.

\begin{lemma}[Equal-bin quadrature]\label{lem:BV}
If a probability $\xi$ assigns mass $1/N$ to each of the $N$ equal bins
of $[0,1]$, then for every function $F$ of bounded variation,
\begin{equation}\label{eq:BV}
 \left|\int F\,d\xi-\int F\,dU\right|\leq\frac{\Var(F)}N.
\end{equation}
\end{lemma}
\begin{proof}
Within each bin the difference between the two conditional averages is
bounded by the oscillation of $F$. Summing with weight $1/N$ proves the
claim, since the sum of the bin oscillations is at most total variation.
At jump endpoints use one-sided representatives consistent with the bin
convention. The template laws~\eqref{eq:template} used below are absolutely
continuous within bins, so endpoint values do not affect their integrals.
\end{proof}

Let $Q=[q_1,q_1+h]\times[q_2,q_2+h]$, $0<h\leq1$, satisfy
\eqref{eq:separation} with respect to the pin $p$. Let $U_Q$ be uniform
probability on $Q$, and let $\rho_Q$ be the affine image of the product
in~\eqref{eq:template}, allowing any integers $A,B\geq1$ and widths
$0<w\leq A^{-1}$, $0<v\leq B^{-1}$. Write
\[
 M_\varepsilon(p,y)=\ind_{\{|y_1-p_1|\geq\varepsilon\}}.
\]
When a norm of a signed pushforward is written below, it means the norm of
its density. Each constituent measure is positive and unnormalized.

\begin{lemma}[Fixed-band comparison]\label{lem:fixed-band}
Both masked distance pushforwards have densities, and
\begin{equation}\label{eq:fixed-band}
 \norm{(D_p)_*(M_\varepsilon\rho_Q)
       -(D_p)_*(M_\varepsilon U_Q)}_\infty
 \leq\frac{4D_0}{h}
       \left(\frac1{v_0A}+\frac1{\varepsilon Bw}\right).
\end{equation}
The pin and the mask are identical in the two terms.
\end{lemma}
\begin{proof}
Use coordinates $X=q_1+hu$, $Y=q_2+hz$ and the local output
$t=|y-p|/h$. Conversion back to physical radius contributes a final
factor $h^{-1}$ to the density.

First fix a horizontal strip $i/A\leq u\leq i/A+w$, keep the part
$|X-p_1|\geq\varepsilon$, and average in the vertical coordinate.
For a fixed local output $t$, set $r=ht$. Horizontal coarea gives the
$t$-density
\begin{equation}\label{eq:horizontal-coarea}
 G_i(z,t)=\frac1w\sum_{\text{allowed roots }X}\frac{r}{|X-p_1|},
 \qquad X-p_1=\pm\sqrt{r^2-(Y-p_2)^2}.
\end{equation}
Indeed $\partial t/\partial u=(X-p_1)/r$, so no extra factor of $h$
belongs in~\eqref{eq:horizontal-coarea}. On each branch, $Y-p_2>0$
and $|X-p_1|$ decreases as $z$ increases. The reciprocal derivative
$r/|X-p_1|$ is monotone. Requiring the root to lie in the strip cuts out
an interval of $z$, and imposing the band cutoff only shortens it.
The zero extension of each monotone interval function has variation at
most twice its supremum. There are at most two branches, each at most
$D_0/(\varepsilon w)$, so
\begin{equation}\label{eq:vertical-variation}
 \Var_z G_i(z,t)\leq\frac{4D_0}{\varepsilon w}.
\end{equation}
Lemma~\ref{lem:BV} replaces $\zeta_{B,v}$ by $U$ at a density cost
$4D_0/(\varepsilon Bw)$. Averaging the strips with their $1/A$ weights
preserves that bound. The width $v$ has not entered except through
$v\leq1/B$.

With the vertical coordinate uniform, fix $u$. There is only one allowed
vertical root because $Y-p_2\geq v_0$. Its local-output density is
\begin{equation}\label{eq:vertical-coarea}
 J(u,t)=\ind_{\{\text{allowed vertical root}\}}
        \ind_{\{|X-p_1|\geq\varepsilon\}}
        \frac{r}{Y-p_2}.
\end{equation}
It is at most $D_0/v_0$. On either side of $X=p_1$ it is monotone in
$|X-p_1|$; requiring the root to lie in the vertical side of $Q$ leaves
at most one interval per side. Therefore
$\Var_u J(u,t)\leq4D_0/v_0$. A second application of
Lemma~\ref{lem:BV}, this time to $\xi_{A,w}$, costs $4D_0/(v_0A)$.
Adding the two errors and returning to physical radius proves the claim.
Coarea is applied on strictly monotone intervals; the cutoff removes the
only horizontal critical point. All bounds hold for almost every output,
which is sufficient for the stated essential supremum.
\end{proof}

\begin{lemma}[Harmonic improvement]\label{lem:harmonic}
At the particular physical cutoff $\varepsilon=h/A$,
\begin{equation}\label{eq:harmonic}
 \norm{(D_p)_*(M_{h/A}\rho_Q)-(D_p)_*(M_{h/A}U_Q)}_\infty
 \leq C D_0\left(\frac1{hv_0A}
              +\frac{\log(2A)}{h^2Bw}\right).
\end{equation}
\end{lemma}
\begin{proof}
For the physical strip
$I_i=[q_1+hi/A,q_1+hi/A+hw]$, let $d_i$ be the infimum of
$|X-p_1|$ over its nonempty retained part. Empty strips are omitted.
Then $d_i\geq h/A$, and the sharper variation bound in the first part
of the preceding proof is $4D_0/(wd_i)$. Each strip lies in its own
coarse bin of length $h/A$. Order bins to the left and right of $p_1$.
Except for at most a fixed number of nearest neighbors, the $l$-th bin
has distance at least $lh/(2A)$. The cutoff controls the nearest bins.
The same ordering works if $p_1$ lies outside $Q$. Hence
\begin{equation}\label{eq:harmonic-sum}
 \frac1A\sum_i d_i^{-1}
 \leq\frac Ch\left(1+\sum_{l=1}^A\frac1l\right)
 \leq\frac Ch\log(2A).
\end{equation}
Average this sharper variation over strips before the vertical quadrature
step. Its local-output error is at most $CD_0\log(2A)/(hBw)$;
physical scaling gives the second term in~\eqref{eq:harmonic}.
The horizontal quadrature term is unchanged.
\end{proof}
The resolving relation $Bw=A\to\infty$ expresses a vertical sampling
spacing fine enough to resolve an actual horizontal interval. A Frostman
bound or a numerical branching profile does not imply either lemma.
